% !TeX root = ../../Book1.tex
% 纲要标题：### B.3 Reed–Solomon 码的代数构造
% 纲要位置：计划纲要.md，第 868 行。
% BEGIN APPENDIX OUTLINE SYNC
% #### B.3.1 求值构造与参数
% #### B.3.2 插值与单错误纠正
% END APPENDIX OUTLINE SYNC

\section{Reed–Solomon 码的代数构造}
\label{b1:sec:B03}

\subsecref{b1:subsec:1704:03}用一次多项式在五个点的取值说明了纠错的可能性。现在把次数界和取值点数作为参数：根数界决定距离，插值负责恢复信息，而错误位置未知时，可以把位置也编码成一个待求多项式。本节只证明一个错误的纠正算法，避免让附录承担完整的解码理论。

\subsection{求值构造与参数}
\label{b1:subsec:B03:01}

\begin{definition}\label{b1:def:B-reed-solomon}
设 \(\mathbb F_q\) 为有限域，\(k,n\) 为满足 \(1\leq k\leq n\leq q\) 的整数，并取两两不同的元素 \(a_1,\ldots,a_n\in\mathbb F_q\)。所有次数小于 \(k\) 的多项式（包含零多项式）在这些点的取值构成线性码
\[
C=\Bigl\{\,\bigl(f(a_1),\ldots,f(a_n)\bigr)\mid
f\in\mathbb F_q[x],\ \deg f<k\,\Bigr\}
\]
称为此处采用的\term[Reed-Solomon-ma]{Reed–Solomon 码}[Reed–Solomon code]。\footnote{里德（Irving Stoy Reed，1923-2012），美国数学家与工程师，研究纠错码，与 Solomon 共同提出 Reed--Solomon 码。}\footnote{所罗门（Gustave Solomon，1930-1996），美国数学家与工程师，研究代数编码，与 Reed 共同提出 Reed--Solomon 码。}条件 \(\deg0<k\) 在该式中按惯例成立。
\end{definition}

条件 \(n\leq q\) 来自取值点必须两两不同：\(\mathbb F_q\) 中一共只有 \(q\) 个可选元素。这一名称在文献中有关于取值点的不同约定；本附录允许任意互异有限域元素，不另加非零坐标乘子，也不加入无穷远点。使用其他版本的参数结论时，应先对照定义。

\begin{theorem}\label{b1:thm:B-rs-parameters}
设 \(\mathbb F_q\) 为有限域，\(k,n\) 为满足 \(1\leq k\leq n\leq q\) 的整数，\(a_1,\ldots,a_n\in\mathbb F_q\) 两两不同，并令
\[
C=\{\,\bigl(f(a_1),\ldots,f(a_n)\bigr)\mid
f\in\mathbb F_q[x],\ \deg f<k\,\}.
\]
则 \(C\) 是参数为 \([n,k,n-k+1]\) 的 Reed--Solomon 码，并达到任意非零 \([n,k,d]\) 线性码都满足的界 \(d\leq n-k+1\)。
\end{theorem}

\begin{proof}
求值映射保持加法与标量乘法。若非零多项式 \(f\) 的次数小于 \(k\)，则 \(\deg f\leq k-1<n\)，由\cref{b1:prop:factor-root}的根数界，它不可能在全部 \(n\) 个不同取值点为零。因此求值映射的核为零，像的维数为 \(k\)。非零 \(f\) 至多在 \(k-1\) 个取值点为零，因此\eqref{b1:eq:B-linear-distance}给出 \(d(C)\geq n-k+1\)。反过来，取
\[
f(x)=\prod_{i=1}^{k-1}(x-a_i).
\]
当 \(k=1\) 时乘积取一。这个多项式恰在前 \(k-1\) 个取值点为零，其码字重量为 \(n-k+1\)，从而取得等号。

最后设 \(D\subseteq\mathbb F_q^n\) 是任意非零 \([n,k,d]\) 线性码，任取 \(d-1\) 个坐标，并固定删去它们。所得映射在 \(D\) 上单射，否则两个不同码字只在被删去的位置不同，距离至多为 \(d-1\)，与最小距离为 \(d\) 矛盾。论证不依赖所删坐标的选择，所以 \(q^k=|D|\leq q^{n-d+1}\)，即 \(d\leq n-k+1\)。
\end{proof}

设 \(\mathbb F_q\) 为有限域，\(n,k,d\) 为正整数，\(C\subseteq\mathbb F_q^n\) 为非零 \([n,k,d]\) 线性码。不等式 \(d\leq n-k+1\) 称为\term[Singleton-jie]{Singleton 界}[Singleton bound]；\footnote{辛格尔顿（Richard Collom Singleton，1928-2007），美国数学统计学家与计算机科学家，给出了编码理论中的长度、维数与最小距离界。}若 \(d=n-k+1\)，则称 \(C\) 为\term[zuidajuli-kefen-ma]{最大距离可分码}[maximum distance separable code]，简称 MDS 码。这里的“最大”是在固定长度和维数的比较中而言，并不意味着其他编码需求也同时最优。“可分”在此是编码术语，与第十六章讨论的多项式或域扩张的可分性无关。

\ProseParagraphBreak
以系数行向量 \((u_0,\ldots,u_{k-1})\) 表示 \(f=\sum_j u_jx^j\)，一个生成矩阵为
\[
G=\begin{pmatrix}
1&1&\cdots&1\\
a_1&a_2&\cdots&a_n\\
\vdots&\vdots&&\vdots\\
a_1^{k-1}&a_2^{k-1}&\cdots&a_n^{k-1}
\end{pmatrix}.
\]
从任何 \(k\) 个未损坏的坐标，都能由\cref{b1:thm:lagrange-interpolation}唯一恢复 \(f\)：取值点互异，插值次数至多为 \(k-1\)，恰与信息多项式的次数界一致。若损坏位置已经标明，则称这些缺失坐标为\term[cachu]{擦除}[erasure]；最多 \(n-k\) 次擦除可以恢复。这个结论不能解释为可纠正 \(n-k\) 个未知位置的错误：擦除允许我们先剔除已知损坏的位置，错误值却会混在其余数据中，不能直接指定哪些数据可供插值。

\ProseParagraphBreak
具体地，若保留下来的 \(k\) 个位置组成指标集 \(I\)，其值为 \(y_i\)，则恢复式是
\[
f(x)=\sum_{i\in I}y_i\prod_{\substack{j\in I\\j\ne i}}
\frac{x-a_j}{a_i-a_j}.
\]
分母非零正是取值点互异的用途。若错误位置没有标出，同一公式依然能拟合任选的 \(k\) 项数据，却未必给出原信息。例如一次错误恰落在所选位置中，插值就会把错误值也当作必须满足的条件。下面引入的多项式 \(E\)，就是为了让未知错误位置的条件自动失去约束，而保留其余位置的正确等式。

\subsection{插值与单错误纠正}
\label{b1:subsec:B03:02}

\begin{samepage}
设 \(n\geq k+2\)，信息多项式 \(f\in\mathbb F_q[x]\) 满足 \(\deg f<k\)，收到 \(y=(y_1,\ldots,y_n)\)，且已知它与 \(f\) 的取值码字至多有一个坐标不同。若错误发生在取值点 \(a_j\)，取首一一次多项式 \(E=x-a_j\)，称为本算法的\term[cuowei-duoxiangshi]{错误定位多项式}[error locator polynomial]。它在错误位置为零，其余位置满足 \(y_i=f(a_i)\)。令 \(Q=Ef\)，则
\begin{equation}\label{b1:eq:B-single-error-system}
Q(a_i)=y_i\cdot E(a_i)\quad(1\leq i\leq n),\qquad
E=x+e_0,\quad \deg Q\leq k.
\end{equation}
实际求解时错误位置尚未知晓。若同时求信息多项式与错误点，直接展开 \((x-a_j)f\) 会产生未知系数之间的乘积。为此不要求先写出 \(f\)，而把 \(e_0\) 及 \(Q=q_0+q_1x+\cdots+q_kx^k\) 的系数视为独立未知量。\eqref{b1:eq:B-single-error-system}于是成为关于 \(q_0,\ldots,q_k,e_0\) 的线性方程组；已知量 \(y_i\) 只是系数，并未产生未知量之间的乘积。若确有错误，\(e_0=-a_j\)，错误取值点是 \(-e_0\)，而不是 \(e_0\)。
\end{samepage}

\ProseParagraphBreak
每一行的形式为
\[
q_0+a_iq_1+\cdots+a_i^kq_k-y_i e_0=y_i a_i.
\]
因此未知数共有 \(k+2\) 个，可以在 \(\mathbb F_q\) 中用线性方程组的消元法求解。若不固定 \(E\) 的首项系数，\((E,Q)\) 和 \((\lambda E,\lambda Q)\) 对任意非零 \(\lambda\) 都满足相同的求值等式。规定 \(E\) 首一既排除 \(E=Q=0\) 这种不携带信息的解，也从非零倍数中选定一个代表；这并不保证辅助变量本身唯一。方程数至少等于未知数个数本身也不保证解存在或解码正确；下一命题将利用“至多一个错误”和多项式根数界证明这两点。

\begin{proposition}\label{b1:prop:B-single-error-decoding}
设 \(\mathbb F_q\) 为有限域，\(k,n\) 为满足 \(1\leq k\leq n-2\) 的整数，\(a_1,\ldots,a_n\in\mathbb F_q\) 两两不同，\(f\in\mathbb F_q[x]\) 满足 \(\deg f<k\)。若收到的向量 \(y=(y_1,\ldots,y_n)\) 与 \(\bigl(f(a_1),\ldots,f(a_n)\bigr)\) 至多有一个坐标不同，则关于
\[
E=x+e_0,\qquad \deg Q\leq k
\]
的线性方程组
\[
Q(a_i)=y_i\cdot E(a_i)\qquad(1\leq i\leq n)
\]
有解，而且每个解都满足 \(Q=Ef\)。因此先解该线性方程组，再作多项式除法 \(f=Q/E\)，即可纠正至多一个错误。
\end{proposition}

\begin{proof}
恰有一次错误时，取其位置对应的 \(E=x-a_j\) 和 \(Q=Ef\)，就在错误处得到两边为零，在所有正确位置得到同一乘积，故存在解。没有错误时，可任取首一一次多项式 \(E\)，仍以 \(Q=Ef\) 得到解。此时 \(E\) 只是为统一方程形式引入的辅助变量，它的根不表示实际错误位置；不同 \(E\) 给出不同辅助解，但不改变商 \(Q/E=f\)。

对方程组的任意解 \((E,Q)\)，只能在正确位置将 \(y_i\) 换成 \(f(a_i)\)，从而得到 \((Q-Ef)(a_i)=0\)。这样的不同取值点至少有 \(n-1\) 个，而 \(\deg(Q-Ef)\leq k<n-1\)。由\cref{b1:prop:factor-root}，\(Q-Ef\) 只能是零多项式。因此 \(Q=Ef\)，且 \(E\) 首一一次，非零，整环中的除法商唯一。这证明了从任意解都恢复同一个 \(f\)。
\end{proof}

这一步是多项式整除，不能改成在各个取值点除以 \(E(a_i)\)：错误位置上的分母为零。若没有错误数的先验保证，方程组无解时须拒绝输出；有解时也不能仅凭解的存在输出信息，还须检查 \(E\mid Q\)、所得商的次数小于 \(k\)，以及重新编码后的码字与 \(y\) 相差至多一个坐标。不符合任一条件也须拒绝输出。这些核验能够确认收到的向量与输出码字相距至多一；由\cref{b1:thm:B-rs-parameters}，此时最小距离 \(n-k+1\geq3\)，两个这样的候选码字彼此至多相差两个坐标，所以只能相同。要认定这个唯一候选就是实际发送码字，仍须依赖传输错误数至多为一的先验保证。

\ProseParagraphBreak
例如在 \(\mathbb F_7\) 上取 \(a_i=0,1,2,3,4\)、\(k=3\)。信息多项式 \(f=1+2x+3x^2\) 编成 \((1,6,3,6,1)\)。若收到 \(y=(1,6,4,6,1)\)，解方程组可得
\[
E=x+5,\qquad Q=3x^3+3x^2+4x+5.
\]
直接相乘有 \(Q=(x+5)(1+2x+3x^2)\)，所以恢复了 \(f\)。这里 \(e_0=5\)，\(E\) 的根为 \(-5=2\)，正是第三个取值点；具体代入检查放在下一节。这一例子中的 \([5,3,3]\) 参数也由\cref{b1:thm:B-rs-parameters}得到，最小距离三保证一次错误的唯一纠正。
